\begin{frame}{RBF의 무한차원 실체, 실효 차원, 머서 조건}
  ``RBF 커널은 무한차원의 $\phi$ 에 대응한다'' — 구체적으로:
  \begin{itemize}
    \item $\|x-x'\|^2 = \|x\|^2+\|x'\|^2 - 2x\cdot x'$ 로 넣고
      $e^{2\gamma\, x\cdot x'}$ 를 테일러 전개하면, $\phi(x)$ 는
      \textbf{모든 도수의 모든 단항식}에 가우시안 인자를 둔
      \textbf{무한}개의 성분 $\phi_\alpha(x)$ 를 가진다.
    \item \kb{실효 차원은 데이터 개수로 제한된다}: $m$개 데이터의
      커널(Gram) 행렬 $K_{ij}=K(x_i,x_j)$ 의 \textbf{랭크는 최대
      $m$} — $\phi$ 가 무한차원이어도 실제로 ``활용하는'' 특징
      차원은 \textbf{데이터 수를 넘지 않음}.
      (실습: 60개 4차원 RBF 커널 행렬이 rank 60임을 고유값으로 확인)
  \end{itemize}
  \vskip0.4em
  \begin{block}{머서(Mercer) 조건 — 유효 커널의 필요 조건}
    커널 함수가 유효하려면 커널 행렬이 \textbf{모든} 데이터에 대해
    \textbf{양의 준정부호}(모든 고유값 $\ge 0$)여야 한다.
    RBF: 60개 4차원($\gamma=0.5$)에서 최소 고유값
    $5.07\times10^{-3}>0$, 랭크 60 — 직접 확인.
    유한 절단만으로도 꽤 정확한 근사 — ``무한차원''은
    \textbf{개념적} 차원일 뿐 계산 비용이 무한하다는 뜻이 \textbf{아님}.
  \end{block}
\end{frame}
